\chapter{KẾT LUẬN VÀ ĐỀ XUẤT}

\section{Kết luận chung}

Đồ án đã hoàn thành việc xây dựng và đánh giá một hệ thống dự báo nồng độ bụi mịn $PM_{2.5}$ tại Việt Nam, bao phủ toàn bộ quy trình từ thu thập dữ liệu thô đến triển khai ứng dụng web phục vụ tra cứu trực tuyến. Hệ thống được thiết kế theo kiến trúc pipeline gồm 5 khối chức năng liên kết tuần tự (thu thập --- tiền xử lý --- huấn luyện --- đánh giá --- triển khai), vận hành trên dữ liệu từ 12 trạm quan trắc thuộc 2 cụm vùng khí hậu (Bắc 6 trạm, Nam 6 trạm) trong giai đoạn 01/2023 -- 12/2025, tần suất quan trắc 1 bản ghi/giờ. Quy mô thực nghiệm đạt 180 tổ hợp đánh giá (6 mô hình $\times$ 5 tầm nhìn dự báo $\times$ 3 chiến lược chia dữ liệu $\times$ 2 vùng), được đo lường đồng thời trên 4 chỉ số thống kê (MAE, RMSE, $R^2$, MAPE). Các kết luận chính rút ra từ quá trình nghiên cứu được trình bày theo từng mục tiêu đã đề ra.

\subsection{Về thu thập và tiền xử lý dữ liệu}

Pipeline thu thập tự động từ hai nguồn API --- Weatherbit cho dữ liệu chất lượng không khí (7 biến: AQI, $PM_{2.5}$, $PM_{10}$, CO, $NO_2$, $SO_2$, $O_3$) và Open-Meteo cho dữ liệu khí tượng lịch sử ERA5 (17 biến thuộc 7 nhóm) --- đã tổng hợp được khoảng 1,63 triệu bản ghi thô trên 32 trạm. Quy trình sàng lọc 6 tiêu chí thiết kế theo 3 cấp độ (toàn vẹn dữ liệu, quan hệ không gian, đa dạng địa lý) đã giảm từ 32 trạm ban đầu xuống còn 12 trạm đạt chuẩn, tương ứng tỷ lệ 37,5\%. Tỷ lệ loại bỏ 62,5\% phản ánh thực trạng chất lượng dữ liệu quan trắc môi trường tại Việt Nam: tỷ lệ dữ liệu khuyết cao (một số trạm vượt 40\%), hiện tượng đóng băng cảm biến phổ biến, và sự trùng lặp giữa các trạm cùng khu vực (hệ số tương quan $r > 0,99$). Kết quả này cho thấy bước sàng lọc chất lượng là bắt buộc trước khi đưa dữ liệu vào mô hình, đặc biệt khi sử dụng các kiến trúc học sâu vốn nhạy cảm với nhiễu đầu vào.

Giai đoạn làm sạch 7 bước --- gồm 3 giai đoạn ghép nối và chuẩn bị, phát hiện và gán nhãn lỗi, làm sạch và phục hồi --- đã xử lý ba loại lỗi chính: dữ liệu đóng băng (phát hiện bằng độ lệch chuẩn trượt $< 10^{-6}$ trên cửa sổ 12 giờ cho 11 biến), giá trị ngoại lai (ngưỡng vật lý kết hợp $3 \times IQR$), và lỗi cảm biến quang học (đối chiếu chéo 3 điều kiện gồm đột biến $PM_{2.5}$, đột biến $PM_{10}$ và ổn định khí gas). Giá trị khuyết sau phát hiện lỗi được phục hồi bằng nội suy tuyến tính cho khoảng nhỏ ($\le 6$ giờ) và K láng giềng gần nhất ($K = 12$) cho khoảng lớn. Toàn bộ chuỗi kết quả sau làm sạch đạt tính liên tục hoàn toàn trên 25.536 mốc giờ cho mỗi trạm, đủ điều kiện cho các phép tính cửa sổ trượt và xây dựng đặc trưng ở giai đoạn tiếp theo.

\subsection{Về kỹ thuật đặc trưng và chuẩn hóa}

Hệ thống 134 đặc trưng phân thành 8 nhóm được thiết kế theo nguyên tắc phân tầng: từ trạng thái hiện tại (nhóm 1--3 gồm 37 biến thô, tương tác và cờ chất lượng), thông tin lịch sử (nhóm 4--6 gồm 44 biến trễ và thống kê trượt, toàn bộ dịch chuyển 1 bước thời gian chống rò rỉ dữ liệu), tín hiệu chu kỳ (nhóm 7 gồm 14 biến mã hóa thời gian bằng hàm lượng giác $\sin/\cos$), đến dữ liệu khí tượng tương lai (nhóm 8 gồm khoảng 20 biến dịch chuyển theo tầm nhìn dự báo). Trong số đó, 10 đặc trưng tương tác được thiết kế dựa trên nguyên lý hóa khí quyển (tiềm năng oxy hóa $O_3 \times (SO_2 + NO_2)$, nguy cơ sulfate ẩm bằng tích độ ẩm và $SO_2$, ổn định nhiệt bằng hiệu nhiệt độ không khí và nhiệt độ đất) --- đây là các biến mà theo tổng quan tài liệu tại Chương~1, chưa có nghiên cứu nào ứng dụng cho bài toán dự báo $PM_{2.5}$ tại Việt Nam.

Chiến lược chuẩn hóa theo 8 nhóm --- lựa chọn tổ hợp phương pháp (log1p, RobustScaler, StandardScaler, MinMaxScaler, cắt phân vị hoặc không chuẩn hóa) phù hợp với đặc tính phân phối của từng nhóm biến --- đã đảm bảo dữ liệu đầu vào đồng nhất về thang đo mà vẫn bảo toàn đặc trưng phân phối riêng biệt. Các bộ chuẩn hóa chỉ được khớp trên tập huấn luyện và lưu trữ dưới dạng tệp .pkl theo từng trạm, phục vụ biến đổi ngược khi suy luận. Hiệu suất dẫn đầu của XGBoost --- mô hình sử dụng trực tiếp toàn bộ hệ thống 230+ đặc trưng (134 đặc trưng cơ sở cộng thêm lag mở rộng, rolling mở rộng, spatial neighbor và station one-hot) --- trên block~5 và block~7 cho thấy hệ thống đặc trưng thủ công cung cấp thông tin dự báo có giá trị cao, đặc biệt khi dữ liệu huấn luyện liên tục hạn chế.

\subsection{Về kết quả đối chuẩn mô hình}

Phân tích trên 180 cấu hình thí nghiệm dẫn đến ba phát hiện chính mang tính hệ thống.

Phát hiện thứ nhất là không tồn tại mô hình đơn lẻ nào vượt trội tuyệt đối trên mọi cấu hình, và hiệu suất tương đối giữa các kiến trúc phụ thuộc mạnh vào sự kết hợp giữa chiến lược chia dữ liệu và đặc điểm vùng địa lý. XGBoost đạt MAE thấp nhất và hệ số suy giảm (slope) chậm nhất trên block~5 và block~7, coi là lựa chọn tối ưu cho kịch bản dữ liệu huấn luyện liên tục hạn chế (120--168 bước thời gian mỗi khối). Cụ thể, tại block~7 miền Bắc, XGBoost đạt MAE $= 6{,}17$ tại $T{+}1$ và $12{,}87$ tại $T{+}24$, với chênh lệch ($\Delta$MAE) chỉ $6{,}70~\mu g/m^3$ --- thấp hơn đáng kể so với iTransformer ($\Delta = 12{,}05$) và XLinear ($\Delta = 11{,}36$). Tuy nhiên, tại block~30 (720 bước thời gian liên tục mỗi khối), các mô hình học sâu thu hẹp khoảng cách một cách đáng kể: iTransformer cải thiện $R^2$ từ 0,06 (block~7) lên 0,40 (block~30) tại miền Bắc $T{+}24$, tức tăng gấp 6,7 lần; E-STGCN ghi nhận slope thấp nhất toàn bảng (0,21 miền Bắc, 0,30 miền Nam) tại block~30, thấp hơn cả XGBoost (0,26 và 0,54). Tại block~30 miền Bắc $T{+}1$, ST-XLinear (MAE $= 5{,}84$) thậm chí vượt qua XGBoost (MAE $= 6{,}42$). Những kết quả này phản ánh sự đánh đổi giữa yêu cầu dữ liệu và khả năng biểu diễn: mô hình cây quyết định cần ít dữ liệu nhưng bị giới hạn bởi đặc trưng thủ công, trong khi mô hình học sâu cần nhiều dữ liệu nhưng có tiềm năng tự học biểu diễn phức tạp hơn.

Phát hiện thứ hai là chênh lệch hiệu suất giữa hai vùng mang tính hệ thống và không phụ thuộc kiến trúc mô hình. RMSE tại miền Nam luôn cao hơn miền Bắc với tỷ lệ dao động 1,1 lần tại $T{+}1$ đến 2,7--4,5 lần tại $T{+}24$ tùy mô hình. Đồng thời, MAE tại $T{+}1$ miền Nam lại thấp hơn miền Bắc (4,04--8,79 so với 5,84--9,34 trên block~5). Sự mâu thuẫn biểu kiến giữa MAE thấp nhưng RMSE cao ở miền Nam được giải thích bằng đặc tính phân phối: khu vực phía Nam có giá trị nền $PM_{2.5}$ thấp phổ biến (dưới $20~\mu g/m^3$) xen lẫn các đợt ô nhiễm cực đoan vượt $200~\mu g/m^3$, tạo ra phân phối đuôi dài. Phân phối dạng này khiến đa số dự báo có sai lệch nhỏ (MAE thấp), nhưng các sai lệch cực đoan bị phóng đại qua phép bình phương trong RMSE. Ba nguyên nhân gốc rễ bao gồm: sự khác biệt về phân phối nồng độ, đặc trưng khí tượng (miền Bắc có tính mùa vụ ổn định do gió mùa Đông Bắc, miền Nam có biến thiên phi tuyến do chuyển đổi nhanh mùa khô--mưa), và mật độ không gian của mạng lưới trạm (6 trạm miền Bắc tập trung trong đồng bằng sông Hồng với tương quan không gian cao, 6 trạm miền Nam phân bố rải rác từ đồng bằng đến ven biển).

Phát hiện thứ ba liên quan đến ảnh hưởng của chiến lược chia dữ liệu. Chiến lược block-based split tĩnh chia dữ liệu theo khối ngày tuần hoàn, trong đó block~30 cung cấp khoảng 720 bước thời gian liên tục mỗi khối --- gấp 6 lần block~5 (120 bước) và gấp 4,3 lần block~7 (168 bước). Ảnh hưởng khác nhau rõ rệt theo nhóm kiến trúc: XGBoost có MAE dao động chỉ $0{,}20~\mu g/m^3$ qua ba block tại miền Bắc $T{+}24$ (biên độ 12,67--12,87); iTransformer ghi nhận bước nhảy $R^2$ lớn nhất ($+0{,}34$ từ block~7 sang block~30); E-STGCN cải thiện đều đặn nhất qua ba bước nhảy. Hai nhóm kiến trúc phân biệt rõ: nhóm ổn định với mọi block (XGBoost, nhờ không phụ thuộc tính liên tục dữ liệu) và nhóm cải thiện mạnh theo block (iTransformer, E-STGCN, XLinear, ST-XLinear, nhờ cơ chế attention và tích chập đồ thị cần ngưỡng dữ liệu tối thiểu để ước lượng tham số hiệu quả).

\subsection{Về mô hình đề xuất ST-XLinear}

Mô hình ST-XLinear --- kiến trúc đề xuất tích hợp Dynamic Graph Builder vào khung XLinear nhằm khai thác đồng thời phụ thuộc thời gian và không gian --- chưa đạt được kỳ vọng vượt trội so với mô hình nền XLinear trên đa số cấu hình. $R^2$ tại block~30 miền Bắc $T{+}24$ của ST-XLinear (0,37) thấp hơn XLinear (0,40) và E-STGCN (0,41). Tại block~30 miền Nam $T{+}24$, ST-XLinear (MAE~$= 19{,}00$) cũng cao hơn E-STGCN (MAE~$= 15{,}80$) và iTransformer (MAE~$= 15{,}33$).

Kết quả này cho phép rút ra hai bài học kỹ thuật. Bài học thứ nhất, mạng lưới 6 trạm trong mỗi cụm có thể quá thưa để ma trận kề động --- dù tích hợp thông tin hướng gió và vận tốc gió --- tạo ra tín hiệu bổ sung có ý nghĩa so với ma trận kề tĩnh dựa trên khoảng cách. Khi số nút trong đồ thị nhỏ, mỗi nút chỉ có tối đa 5 kết nối khả dụng, và thông tin từ tích chập đồ thị bị giới hạn bởi quy mô lân cận. E-STGCN với ma trận kề tĩnh lại cho kết quả tốt hơn tại block~30, gợi ý rằng quan hệ không gian ổn định có thể phù hợp hơn quan hệ động cho mạng lưới thưa, do cập nhật hướng gió liên tục gây dao động trong quá trình tối ưu khi số mẫu không gian chưa đủ. Bài học thứ hai, cơ chế Attention Pooling tóm tắt đặc trưng thời gian thành Node Summary trước khi đưa vào Spatial GCN có thể gây mất thông tin chi tiết trong chuỗi, đặc biệt ở tầm nhìn dài khi mà chính các chi tiết biến động trong cửa sổ 48 giờ mang tính dự báo cao.

Mặc dù chưa đạt hiệu suất vượt trội, ST-XLinear vẫn cho thấy tính khả thi của kiến trúc lai temporal--spatial: tại block~30 miền Bắc $T{+}1$, ST-XLinear đạt MAE~$= 5{,}84$ (thấp nhất toàn bảng) và MAPE~$= 19{,}49\%$ (thấp nhất toàn bảng). Điều này cho thấy tiềm năng của kiến trúc khi được áp dụng trên mạng lưới dày hơn và với các cải tiến cơ chế tổng hợp không gian phù hợp.

\subsection{Về phương pháp Ensemble}

Phương pháp Ensemble trung bình trọng số tuyến tính --- kết hợp đầu ra XLinear và E-STGCN với trọng số $\alpha$ tối ưu qua Grid Search trên tập xác thực --- cho hiệu quả hạn chế. Chênh lệch $\Delta_{\text{Ensemble}}$ (MAE$_{\text{Ensemble}}$ $-$ MAE$_{\text{min}}$) dương trên 28/30 cấu hình, dao động từ 0,01 đến 3,58~$\mu g/m^3$, tức Ensemble kém hơn hoặc bằng mô hình thành phần tốt nhất trên hầu hết mọi trường hợp.

Phân tích nguyên nhân cho thấy ba yếu tố chính. Thứ nhất, tương quan sai số cao giữa các mô hình thành phần khi được huấn luyện trên cùng bộ dữ liệu và phân chia train/validation/test, khiến chúng mắc lỗi tương tự trên cùng giai đoạn. Thứ hai, chênh lệch hiệu suất lớn giữa mô hình tốt nhất (thường là XGBoost hoặc XLinear) và các mô hình còn lại buộc cơ chế tối ưu trên tập xác thực gán trọng số gần như toàn bộ cho mô hình dẫn đầu, khiến Ensemble suy biến thành mô hình đơn lẻ bị pha loãng. Thứ ba, cơ chế kết hợp tuyến tính giả định sai số đồng nhất về phương sai --- giả định không thỏa mãn khi phân phối $PM_{2.5}$ có đuôi dài, dẫn đến trung bình tuyến tính phóng đại sai lệch phần trăm tại các giá trị nền thấp (MAPE của Ensemble đạt tới 120,80\% tại block~7 miền Nam $T{+}24$, cao nhất trong 6 mô hình).

\subsection{Về chi phí tính toán và triển khai}

Sáu mô hình phân thành ba nhóm chi phí rõ rệt: nhóm chi phí thấp gồm XGBoost (trung bình 72 giây) và E-STGCN (trung bình 183 giây); nhóm chi phí trung bình gồm Ensemble (1.699 giây) và iTransformer (1.879 giây); nhóm chi phí cao gồm ST-XLinear (2.573 giây) và XLinear (3.575 giây). Khoảng cách giữa nhóm thấp nhất và cao nhất lên tới 50 lần (72 giây so với 3.575 giây), đặt ra bài toán đánh đổi giữa hiệu suất dự báo và chi phí vận hành trong triển khai thực tế.

XGBoost đặc biệt phù hợp cho kịch bản triển khai thời gian thực nhờ ba lợi thế: thời gian huấn luyện ngắn nhất (cho phép tái huấn luyện hàng giờ), không yêu cầu GPU, và thời gian suy luận dưới 1 giây cho 5 tầm nhìn dự báo đồng thời. Ứng dụng web trên nền tảng Streamlit tích hợp mô hình XGBoost block~7 đã minh chứng tính khả thi của việc triển khai hệ thống dự báo $PM_{2.5}$ trên phần cứng phổ thông (laptop i5-12500H, 16 GB RAM, GPU RTX 3050) phục vụ tra cứu trực tuyến tại 12 trạm.

\subsection{Tổng hợp đánh giá}

Xét tổng thể 5 mục tiêu ban đầu, 3 mục tiêu đã hoàn thành đạt yêu cầu (pipeline dữ liệu, benchmark 6 mô hình, triển khai ứng dụng web). Mục tiêu kỹ thuật đặc trưng đã hoàn thành ở mức triển khai nhưng chưa có thực nghiệm phân lập để đo lường riêng biệt đóng góp của từng nhóm đặc trưng đối với hiệu suất mô hình. Mục tiêu đề xuất Dynamic Graph (ST-XLinear) đã hoàn thành ở mức kiến trúc và thực nghiệm nhưng chưa đạt hiệu suất vượt trội như kỳ vọng, đồng thời mở ra các giả thuyết cần kiểm chứng trên mạng lưới trạm dày hơn.

Đóng góp chính của đồ án nằm ở ba khía cạnh. Thứ nhất, thiết lập khung đối chuẩn toàn diện nhất từ trước đến nay cho bài toán dự báo $PM_{2.5}$ tại Việt Nam, bao phủ đồng thời Machine Learning cổ điển, Deep Learning temporal, Transformer, và Spatio-Temporal Graph trên cùng bộ dữ liệu quan trắc đa năm, đa vùng, đa horizon và đa chiến lược phân chia --- lấp đầy khoảng trống ``thiếu hụt hệ thống đánh giá đối chuẩn đồng thời'' đã nhận diện tại Chương~1. Thứ hai, xây dựng hệ thống 134 đặc trưng tích hợp kiến thức hóa khí quyển vào quá trình kỹ nghệ đặc trưng cho dữ liệu Việt Nam, khắc phục phần nào khoảng trống ``kỹ nghệ hóa đặc trưng chưa dựa trên nguyên lý khí quyển học''. Thứ ba, phân tích định lượng ảnh hưởng của chiến lược chia dữ liệu đến hiệu suất từng kiến trúc, cung cấp hướng dẫn thực tiễn cho việc lựa chọn mô hình phù hợp với quy mô dữ liệu sẵn có.

\section{Hạn chế}

Mặc dù đã đạt được các kết quả nhất định, nghiên cứu tồn tại một số hạn chế cần được ghi nhận nhằm đặt kết quả trong bối cảnh đúng đắn.

\subsection{Quy mô mạng lưới trạm}

Mạng lưới 12 trạm (6 trạm/cụm) là quy mô nhỏ so với diện tích lãnh thổ. Mật độ trạm thưa khiến các mô hình không gian (E-STGCN, ST-XLinear) không phát huy được tối đa tiềm năng kiến trúc: tích chập đồ thị trên 6 nút chỉ khai thác được tương quan cục bộ giới hạn. Kết quả về hiệu quả (hoặc thiếu hiệu quả) của các kiến trúc không gian chưa thể tổng quát hóa cho mạng lưới dày hơn, nơi ma trận kề phản ánh trung thực hơn cấu trúc vận chuyển khí quyển. Ngoài ra, 12 trạm chỉ bao phủ 13 tỉnh/thành phố, chưa đại diện cho các vùng khí hậu đa dạng khác như Tây Nguyên, Duyên hải miền Trung hay Đồng bằng sông Cửu Long mở rộng.

\subsection{Thiếu thực nghiệm phân lập đặc trưng}

Hệ thống 134 đặc trưng trong 8 nhóm được sử dụng đồng thời cho mô hình, nhưng chưa có thực nghiệm phân lập đo lường đóng góp riêng biệt của từng nhóm. Cụ thể, chưa thể xác định mức độ đóng góp biên của nhóm đặc trưng tương tác hóa khí quyển (nhóm 2) so với các nhóm trễ truyền thống (nhóm 4--6), hoặc mức ảnh hưởng của nhóm thời tiết tương lai (nhóm 8) đối với từng tầm nhìn dự báo cụ thể. Việc thiếu thông tin này khiến khó xác định hướng tối ưu hệ thống đặc trưng cho các ứng dụng trong tương lai.

\subsection{Phương pháp Ensemble đơn giản}

Phương pháp trung bình trọng số tuyến tính là cơ chế kết hợp đơn giản nhất. Nghiên cứu chưa khảo sát các phương pháp kết hợp phi tuyến hoặc kết hợp thích ứng theo miền giá trị. Kết luận ``Ensemble không hiệu quả'' cần được giới hạn cho phương pháp trung bình trọng số tuyến tính, không nên mở rộng cho toàn bộ lớp phương pháp ensemble learning.

\subsection{Đánh giá chưa kèm khoảng tin cậy}

Phương pháp block-based split tạo ra một bộ kết quả duy nhất cho mỗi cấu hình thay vì phân phối kết quả. Các chênh lệch hiệu suất nhỏ giữa mô hình (ví dụ: MAE 5,84 của ST-XLinear so với 5,88 của XLinear tại block~30 miền Bắc $T{+}1$) chưa thể khẳng định có ý nghĩa thống kê hay chỉ là biến động do cách phân chia dữ liệu. Phương pháp xác thực chéo chuỗi thời gian với nhiều lần lặp sẽ cho phép kiểm định thống kê các chênh lệch này.

\subsection{Thiếu đường cơ sở bên ngoài}

Nghiên cứu chỉ so sánh nội bộ giữa 6 kiến trúc Machine Learning/Deep Learning mà chưa đối chiếu với mô hình mô phỏng vật lý--hóa học --- vốn đang là công cụ chủ đạo trong các nghiên cứu dự báo chất lượng không khí tại Việt Nam. Thiếu vắng đường cơ sở này khiến khó xác định liệu nhóm phương pháp thống kê/học máy có thể thay thế hoặc bổ trợ cho mô phỏng số ở quy mô khu vực.

\section{Kiến nghị và hướng phát triển}

Dựa trên các kết quả đạt được và hạn chế đã nhận diện, phần này trình bày các kiến nghị cải tiến và hướng phát triển thang hệ thống dự báo cho các nghiên cứu tiếp theo. Mỗi hướng phát triển được liên kết trực tiếp với một hoặc nhiều hạn chế cụ thể, nhằm đảm bảo tính khả thi và tính hệ thống.

\subsection{Mở rộng quy mô mạng lưới trạm quan trắc}

Mạng lưới 12 trạm (6 trạm mỗi cụm) hiện tại chưa đủ mật độ để các mô hình không gian phát huy tối đa tiềm năng kiến trúc. Kết quả thực nghiệm cho thấy E-STGCN đạt slope thấp nhất tại block~30 (0,21 miền Bắc) nhờ khai thác tương quan liên trạm qua ma trận kề tĩnh Gaussian kernel, trong khi ST-XLinear với Dynamic Graph chưa tạo ra giá trị bổ sung rõ rệt trên mạng lưới thưa --- $R^2$ tại block~30 miền Bắc $T{+}24$ của ST-XLinear (0,37) thấp hơn cả E-STGCN (0,41) lẫn XLinear nền (0,40).

Việc mở rộng lên 30--50 trạm, đặc biệt gia tăng mật độ tại các khu vực đô thị lớn (Hà Nội, TP. Hồ Chí Minh), các cụm công nghiệp trọng điểm (Bắc Ninh, Bình Dương, Đồng Nai) và các vùng khí hậu chưa được bao phủ (Tây Nguyên, Duyên hải miền Trung), sẽ mang lại hai lợi ích trực tiếp. Thứ nhất, mạng lưới dày hơn cung cấp đủ ngữ cảnh không gian để kiểm chứng giả thuyết thiết kế cốt lõi của Dynamic Graph Builder: liệu ma trận kề cập nhật liên tục theo hướng gió có thực sự phản ánh trung thực hơn quá trình vận chuyển ô nhiễm so với ma trận kề tĩnh, khi mỗi nút có 10--20 kết nối thay vì chỉ 5 như hiện tại. Thứ hai, mạng lưới đa vùng bao phủ nhiều chế độ khí hậu khác nhau cho phép đánh giá tính tổng quát hóa của mô hình vượt ra ngoài hai cụm Bắc--Nam hiện tại, từ đó xác định liệu các mẫu hình hiệu suất quan sát được (chênh lệch vùng, ảnh hưởng block) có bền vững hay phụ thuộc vào đặc tính cụ thể của dữ liệu đồng bằng sông Hồng và đồng bằng sông Cửu Long.

Về mặt kỹ thuật, mở rộng mạng lưới đặt ra thách thức về phép tích chập đồ thị trên đồ thị lớn: ma trận kề kích thước $50 \times 50$ làm tăng chi phí tính toán và bộ nhớ gấp khoảng 70 lần so với ma trận $6 \times 6$ hiện tại. Các phương pháp lấy mẫu đồ thị  hoặc tích chập đồ thị thưa cần được xem xét để duy trì tính khả thi tính toán.

\subsection{Thực nghiệm phân lập đặc trưng}

Hệ thống 134 đặc trưng trong 8 nhóm hiện được sử dụng đồng thời mà chưa có đánh giá định lượng đóng góp biên của từng nhóm. Để xác định hướng tối ưu, nghiên cứu tiếp theo cần thực hiện thực nghiệm phân lập theo hai cách tiếp cận bổ sung. Cách tiếp cận thứ nhất là loại bỏ tuần tự: huấn luyện mô hình XGBoost 8 lần, mỗi lần loại bỏ 1 trong 8 nhóm đặc trưng, đo mức suy giảm MAE tương ứng. Nhóm gây suy giảm lớn nhất khi bị loại bỏ là nhóm có đóng góp biên cao nhất. Cách tiếp cận thứ hai là thêm dần: bắt đầu từ nhóm đặc trưng thô (nhóm 1), lần lượt thêm từng nhóm theo thứ tự và đo mức cải thiện tích lũy. Hai cách tiếp cận cho kết quả bổ sung: loại bỏ tuần tự đánh giá đóng góp trong ngữ cảnh đầy đủ, thêm dần đánh giá đóng góp tăng dần.

Đặc biệt, cần đo lường riêng ảnh hưởng của nhóm đặc trưng tương tác hóa khí quyển (nhóm 2 gồm tiềm năng oxy hóa, nguy cơ sulfate ẩm, ổn định nhiệt) so với nhóm trễ truyền thống (nhóm 4--6). Nếu nhóm 2 cho đóng góp biên đáng kể, điều này xác nhận giá trị của việc tích hợp kiến thức lĩnh vực vào kỹ nghệ đặc trưng. Nếu đóng góp không đáng kể, nhóm này có thể được loại bỏ để giảm chiều đầu vào và tăng tốc huấn luyện.

\subsection{Cải tiến phương pháp kết hợp mô hình}

Phương pháp Ensemble trung bình trọng số tuyến tính cho thấy hiệu quả hạn chế do ba yếu tố đã phân tích (tương quan sai số, chênh lệch hiệu suất, giả định tuyến tính). Các hướng cải tiến cần tập trung vào việc giải quyết từng yếu tố.

Để giảm tương quan sai số, nghiên cứu có thể áp dụng chiến lược đa dạng hóa dữ liệu huấn luyện: thay vì huấn luyện mọi mô hình trên cùng bộ dữ liệu, mỗi mô hình thành phần sử dụng tập con đặc trưng khác nhau --- một mô hình chỉ sử dụng đặc trưng khí tượng, mô hình khác chỉ sử dụng đặc trưng hóa học, mô hình thứ ba sử dụng toàn bộ. Sự đa dạng về đầu vào dẫn đến đa dạng về mẫu hình sai số, tạo điều kiện thuận lợi cho kết hợp triệt tiêu.

Để khắc phục giả định tuyến tính, phương pháp stacking nên được xem xét: sử dụng mô hình bậc hai (có thể là gradient boosting hoặc mạng nơ-ron nhỏ) để học hàm kết hợp phi tuyến từ các đầu ra mô hình thành phần. Cơ chế kết hợp thích ứng theo miền giá trị cũng khả thi: gán trọng số khác nhau tùy theo dải $PM_{2.5}$ mà mô hình đang dự báo (ví dụ: ưu tiên XGBoost khi $PM_{2.5}$ dự báo thuộc dải thấp, ưu tiên E-STGCN khi thuộc dải cao). Phương pháp Mixture of Experts, trong đó một mạng gating nhỏ quyết định trọng số kết hợp dựa trên ngữ cảnh đầu vào (giờ trong ngày, mùa, điều kiện gió), có thể khai thác sự chuyên biệt hóa theo kịch bản của từng mô hình.

\subsection{Tích hợp dữ liệu dự báo khí tượng số}

Nhóm 8 trong hệ thống đặc trưng (thời tiết tương lai) hiện sử dụng giá trị thực tế quan trắc tại thời điểm tương lai trong pha huấn luyện và đánh giá, nhằm đo lường giới hạn trên của hiệu năng dự báo khi thông tin khí tượng hoàn hảo. Khi triển khai thực tế, giá trị này phải được thay thế bằng đầu ra từ mô hình dự báo khí tượng số như GFS (Global Forecast System) hoặc ECMWF IFS. Sai số dự báo NWP --- vốn tăng theo tầm nhìn dự báo --- sẽ lan truyền vào mô hình $PM_{2.5}$ và gây suy giảm hiệu suất, đặc biệt ở $T{+}12$ và $T{+}24$ giờ.

Để giảm thiểu tác động này, hai chiến lược cần được khảo sát. Chiến lược thứ nhất là huấn luyện mô hình trực tiếp trên đầu ra NWP (thay vì giá trị quan trắc thực tế) cho nhóm 8: điều này cho phép mô hình $PM_{2.5}$ tự học cách bù trừ sai số hệ thống của NWP trong quá trình tối ưu hàm mất mát. Chiến lược thứ hai là bổ sung biến ``sai số NWP'' (hiệu giữa dự báo NWP và giá trị quan trắc tại cùng thời điểm) làm đặc trưng bổ sung, cung cấp cho mô hình thông tin về mức độ tin cậy của dữ liệu khí tượng đầu vào. Cần thực hiện thí nghiệm đối sánh 3 cấu hình (thời tiết thực, thời tiết NWP, thời tiết NWP có hiệu chỉnh) trên cùng mô hình XGBoost để đo lường định lượng mức suy giảm và hiệu quả bù trừ.

\subsection{Áp dụng xác thực chéo chuỗi thời gian}

Phương pháp block-based split hiện tại với 3 cấu hình (block~5, block~7, block~30) đã cung cấp một mức đánh giá tính ổn định qua các chiến lược phân chia, nhưng mỗi cấu hình chỉ tạo ra một bộ kết quả duy nhất. Các chênh lệch hiệu suất nhỏ --- ví dụ MAE~$= 5{,}84$ (ST-XLinear) so với $5{,}88$ (XLinear) tại block~30 miền Bắc $T{+}1$, chênh lệch chỉ $0{,}04~\mu g/m^3$ --- chưa thể khẳng định có ý nghĩa thống kê.

Phương pháp xác thực chéo mở rộng (expanding window cross-validation) phù hợp cho dữ liệu chuỗi thời gian: cửa sổ huấn luyện mở rộng dần qua thời gian, trong khi cửa sổ kiểm thử trượt sang giai đoạn tiếp theo. Với 3 năm dữ liệu (01/2023 -- 12/2025), có thể thiết kế 6--12 lần lặp với cửa sổ kiểm thử 3--6 tháng, tạo ra phân phối kết quả cho mỗi chỉ số đánh giá. Phân phối này cho phép: (i) tính khoảng tin cậy 95\% cho MAE, RMSE, $R^2$ của mỗi mô hình; (ii) thực hiện kiểm định Wilcoxon signed-rank hoặc Friedman test để xác định liệu chênh lệch giữa các mô hình có ý nghĩa thống kê; và (iii) phát hiện liệu các mẫu hình hiệu suất (sự vượt trội của E-STGCN tại block~30, bước nhảy $R^2$ của iTransformer) có bền vững qua nhiều lần phân chia hay chỉ là sản phẩm của đặc điểm giai đoạn thời gian cụ thể.

\subsection{Mở rộng so sánh với đường cơ sở bổ sung}

Nghiên cứu hiện tại tập trung so sánh nội bộ giữa 6 kiến trúc Machine Learning và Deep Learning. Để định vị rõ hơn vị trí của nhóm phương pháp này trong bài toán dự báo chất lượng không khí tại Việt Nam, hai nhóm đường cơ sở cần được bổ sung.

Nhóm thứ nhất là mô hình mô phỏng vật lý --- hóa học, đặc biệt hệ thống CMAQ kết hợp WRF --- vốn đang là công cụ chủ đạo trong các nghiên cứu trong nước. Nguyen et al. (2024) báo cáo mô hình CMAQ kết hợp WRF đạt MAPE xấp xỉ 63\% tại điều kiện phức tạp, cao hơn đáng kể so với XGBoost (21--54\% tùy cấu hình) trong nghiên cứu này. Tuy nhiên, so sánh trực tiếp giữa hai phương pháp trên cùng tập dữ liệu, cùng khoảng thời gian và cùng vị trí trạm vẫn chưa được thực hiện. Thí nghiệm đối sánh này sẽ xác định liệu các phương pháp học máy có thể thay thế (khi dữ liệu quan trắc đủ lớn) hoặc bổ trợ (đầu ra CMAQ làm đặc trưng đầu vào bổ sung) cho mô phỏng số ở quy mô khu vực.

Nhóm thứ hai là các mô hình nền tảng cho chuỗi thời gian --- ví dụ: TimesFM (Google), Chronos (Amazon), Lag-Llama hoặc các biến thể mô hình ngôn ngữ lớn được tinh chỉnh cho dự báo. Các mô hình này được huấn luyện trước trên hàng tỷ điểm dữ liệu chuỗi thời gian đa ngành, và có tiềm năng đạt hiệu suất cao ngay cả khi dữ liệu miền ứng dụng hạn chế --- chính là điều kiện thường gặp trong dữ liệu quan trắc môi trường tại Việt Nam. Việc đánh giá liệu khả năng học chuyển giao ) từ dữ liệu đa miền có mang lại lợi thế cạnh tranh so với các mô hình huấn luyện từ đầu sẽ cung cấp định hướng dài hạn cho hệ thống dự báo.

\subsection{Nâng cấp ứng dụng triển khai}

Ứng dụng web hiện tại hoạt động ở chế độ tra cứu thủ công: người dùng chủ động truy cập giao diện, chọn trạm và nhận kết quả dự báo. Mô hình triển khai này phù hợp cho giai đoạn nghiên cứu và minh chứng tính khả thi (proof of concept), nhưng chưa đáp ứng yêu cầu vận hành trong thực tế khi cần giám sát liên tục và cảnh báo chủ động.

Để nâng cấp hệ thống lên mức vận hành, ba thành phần cần được bổ sung. Thành phần thứ nhất là bộ thu thập dữ liệu tự động hoạt động theo lịch trình cron job với tần suất hàng giờ, đồng bộ dữ liệu mới nhất từ Weatherbit API và Open-Meteo API, thực hiện làm sạch và xây dựng đặc trưng theo cùng pipeline đã thiết kế tại Chương~2, sau đó lưu trữ vào cơ sở dữ liệu trung gian phục vụ suy luận. Thành phần thứ hai là cơ chế tái huấn luyện mô hình định kỳ với tần suất hàng tuần hoặc hàng tháng, nhằm thích ứng với xu hướng phân phối dữ liệu thay đổi theo thời gian --- đặc biệt quan trọng khi chuyển mùa hoặc khi có thay đổi nguồn phát thải trong khu vực. Thời gian huấn luyện XGBoost trung bình 72 giây cho thấy tái huấn luyện định kỳ hoàn toàn khả thi trên phần cứng phổ thông. Thành phần thứ ba là module cảnh báo tự động: khi giá trị $PM_{2.5}$ dự báo tại bất kỳ mốc thời gian nào ($T{+}1$ đến $T{+}24$) vượt ngưỡng cảnh báo theo quy chuẩn (ví dụ: trên $75~\mu g/m^3$ tương ứng mức ``Không tốt cho nhóm nhạy cảm'' theo thang AQI), hệ thống gửi thông báo đẩy qua các kênh tích hợp (email, ứng dụng nhắn tin) đến người dùng đã đăng ký, phục vụ trực tiếp nhu cầu bảo vệ sức khỏe cộng đồng --- đặc biệt trong các đợt ô nhiễm nghiêm trọng vào mùa đông tại khu vực đồng bằng sông Hồng nơi nồng độ $PM_{2.5}$ có thể đạt 50--84~$\mu g/m^3$ ở tầng thấp.

Ngoài ba thành phần trên, giao diện ứng dụng có thể được mở rộng thêm tính năng hiển thị bản đồ nhiệt nội suy không gian --- sử dụng phương pháp Kriging hoặc IDW (Inverse Distance Weighting) trên đầu ra dự báo của 12 trạm để ước lượng nồng độ $PM_{2.5}$ tại các vị trí không có trạm quan trắc. Tính năng này mang lại giá trị trực quan cho người dùng phổ thông và hỗ trợ cơ quan quản lý đánh giá diện rộng chất lượng không khí mà không cần đầu tư thêm hạ tầng trạm quan trắc.

\subsection{Cải tiến kiến trúc Dynamic Graph}

Kết quả thí nghiệm cho thấy ST-XLinear chưa vượt trội như kỳ vọng, tuy nhiên kiến trúc lai temporal--spatial vẫn đạt hiệu suất dẫn đầu tại một số cấu hình cụ thể (MAE thấp nhất tại block~30 miền Bắc $T{+}1$). Để phát huy tiềm năng, ba hướng cải tiến kiến trúc cần được xem xét.

Hướng thứ nhất là thay thế cơ chế Attention Pooling (hiện tóm tắt toàn bộ chuỗi thời gian thành một vector Node Summary duy nhất trước khi đưa vào Spatial GCN) bằng cơ chế tích chập không gian đa bước thời gian: thực hiện tích chập đồ thị tại mỗi bước thời gian $t$ trong cửa sổ 48 giờ, sau đó mới tổng hợp theo chiều thời gian. Cách tiếp cận này bảo toàn thông tin chi tiết về biến động trong cửa sổ mà Attention Pooling có thể đã bỏ sót, với đổi lấy chi phí tính toán tăng tuyến tính theo độ dài cửa sổ.

Hướng thứ hai là áp dụng ma trận kề đa tầng: kết hợp đồng thời ma trận kề tĩnh (khoảng cách Haversine), ma trận kề động theo gió (Dynamic Graph hiện tại), và ma trận kề có thể học --- trong đó ma trận có thể học được tối ưu trực tiếp qua lan truyền ngược, cho phép mô hình tự phát hiện các mối quan hệ không gian ẩn mà khoảng cách vật lý và hướng gió chưa nắm bắt được. Tổng hợp ba nguồn thông tin không gian này có thể giúp mô hình thích ứng với nhiều chế độ vận chuyển ô nhiễm khác nhau.

Hướng thứ ba là tích hợp cơ chế thích ứng trọng số $\alpha_{\text{graph}}$ (hiện cố định, kiểm soát tỷ lệ giữa khoảng cách tĩnh và hướng gió động trong ma trận kề) thành tham số có thể học hoặc thay đổi theo thời gian. Giá trị $\alpha$ tối ưu có thể khác nhau giữa mùa đông (khi nghịch nhiệt mạnh, gió yếu, ảnh hưởng không gian giảm) và mùa hè (khi gió mạnh, lan truyền ô nhiễm qua khoảng cách lớn hơn). Cơ chế thích ứng cho phép mô hình tự điều chỉnh mức độ sử dụng thông tin không gian tùy theo ngữ cảnh khí tượng, thay vì áp dụng hằng số cố định.
